\documentclass[10pt]{article}

\usepackage[utf8]{inputenc}

\usepackage[T1]{fontenc}

\usepackage{amsmath}

\usepackage{amsfonts}

\usepackage{amssymb}

\usepackage[version=4]{mhchem}

\usepackage{stmaryrd}

\usepackage{graphicx}

\usepackage[export]{adjustbox}

\graphicspath{ {./images/} }

\usepackage{mathrsfs}



\DeclareUnicodeCharacter{0131}{$\imath$}



\begin{document}

ФАЗОВЫЙ ПЕРЕХОД - физическое явление, происходящее в макроскопич. системах и состоящее в том, что в нек-рых состояниях равновесия системы сколь угодно малое воздействие приводит к резкому изменению еө свойств: система из одной своей однородной фазы переходит в другую. Математически Ф. п. трактуется как резкое нарушение структуры и свойств т. н. гиббсовских распредөлений, описывающих равнодесные состояния системы, при сколь угодно малых изменениях параметров, определяющих равновесие.



Лит.: [1] ЈI а н д а у Ј. Д., Л и ф шा и ц $\mathbf{E}$ М , Статистическая физика, 3 изд., ч. 1, М., 1976; [2] С и н а й Я.' Г., Теория фазовых переходов, М., 1980.,



ФАЗОВЫХ ИНТЕГРАЛОВ МЕТОД - то жк, что BKЬ-метод.



ФАКТОР - инволютивная подалгебра ㄱ алгебры $B\left(H_{1} H\right)$ линейных операторов в гильбертовом п ространстве $H$, замкнутая относительно т. н. с л а б о й с ходимости олераторов и обладающая тем свойством, что центр (т. е. совокупность всех оператаров из $\mathfrak{U}$, коммутирующих с любым оператором из $\mathfrak{l}$ ) состоит из скалярнократных единичному оператору.



Если $\mathfrak{A}$ есть Ф., то для большого запаса подшространств $F$ из $H$ можно определить понятие размер-



\includegraphics[max width=\textwidth, center]{2023_07_09_b0c21899c65d7aea98ddg-1(1)}

няющегося не при произвольных изометриях $\mathscr{F}$, a только при изометриях из данного Ф. и обладающих дополнительными естествснными свойствами (напр., $\left.\operatorname{dim}_{\Re}\left(F_{1} \oplus F_{2}\right)=\operatorname{dim}_{\vartheta} F_{1}+\operatorname{dim}_{\Im} F_{2}\right)$. Все Ф. разбиваются на пять классов в соответствии со значениями, к-рые может принимать $\operatorname{dim}_{\mathfrak{A}} F$, притем, напр., для Ф. класca $I I_{\infty}$ еe значениями могут быть любые числа из $[0, \infty]$. $\quad$ М. И. Войцеховский.



\includegraphics[max width=\textwidth, center]{2023_07_09_b0c21899c65d7aea98ddg-1(2)}

н о м у д е ли т е л оо $N$ - групша, образуемая смежныли классами $N g, g \in G$, группы $G$ и обозначаемая $G / N$ (см. Нормальный делитель). Умножение смежных нлассов производится по формуле



\[

N g_{1} \cdot N g_{2}==N g_{1} g_{2} .

\]



Единицей Ф. является класс $N=N \cdot 1$, обратным к классу $N g$ класс $N g^{-1}$.



Отображение $x: g \rightarrow N g$ будет эпиморфизмом гругшы $G$ на Ф. $G / N$, наз. К а н о ни ч е с $\kappa$ и м м о р и з м о м $G$ на $G / N$. Если $\varphi: G \rightarrow G^{\prime}-$ произвольный эпиморфизм группы $G$ на группу $G^{\prime}$, то ядро $K$ эıиморфизма $\varphi$ - нормальный делитель груіпы $G$, а факторгруппа $G / K$ изоморфиа групше $G^{\prime}$, точнее, существует изоморфизм $\psi$ группы $G / K$ на группу $G^{\prime}$ такоі, что диаграмма



\begin{center}

\includegraphics[max width=\textwidth]{2023_07_09_b0c21899c65d7aea98ddg-1(3)}

\end{center}



коммутативна, где $x \rightarrow$ канонич. эпиморфиязм $G \rightarrow$ $\rightarrow G / K$.



Ф. групшы $G$ можно определять, исходя из нек-рої конгруәнции на $G$, как множество классов конгруәнтных әлементов относительно умножения классов, Всовозможные конгруәнции на группе находятся во взаимно однозначном соответствии с нормальными делителями группы, а Ф. по конгруәнциям совпадают с Ф. шо нормальным делителям. Ф. является нормальным факторобъектом в категории групп. Н. Н. Вильлмс.



ФАКТОРИАЛ - функция, определенная на множестве целых неотрицательных чисел, значение к-рой равно произведению натуральных чисел от 1 до натурального числа $n$, т. е. $1 \cdot 2 \ldots n$; обозначается $n$ ! (по определению, $0 !=1$ ). При больших $n$ приближенное выражение Ф. дается Стирлинга борлулой. Ф. равен числу перестановок из $n$ алементов. Ф. наз. также более общее выражение



\[

(a)_{\mu}=a(a+1)(a+2) \ldots(a+\mu-1),

\]



где $a$ - любое томплексное число, $\mu$ - натуральное, $(a)_{0}$. См. также Галма-функция. БСЭ-з,



ФАКТОРИАЛЬНОЕ КОЈЦЦО - кольцо с однозначным равложением на множители. Точнее, Ф. к. $A-$ әто область целостности, в к-рой можно выбрать систему әкстремальных элементов $P$ такую, что любой ненулевой әлемент $a \in A$ донускает өдинственное представление вида



\[

a=u \prod_{p \in P} p^{n(p)},

\]



гце $u$ обратим, а целые неотрицательные покаватели $n(p)$ отличны от нуля только для конечного числа әлементов $p \in P$. При этом элемент наз. ә к с т р eм а л ь ны м в $A$, если из $p=u v$ следует, что либо $u$, либо $v$ обратим в $A$, п $p$ необратим в $A$.



В Ф. к. существует наибольший общий делитель и наименьшее общее кратное любых двух элементов. Кольцо $A$ факториально тогда и только тогда, когда оно является кольцом Крулля и выполняется одно из следующи $\mathbf{x}$ эквивалентных условий: (1) любой дивизориальный идеал в $A$ является главным; (2) любой простой идеал высоты 1 главный; (3) любое непустое семейство главных идеалов обладает максимальным элементом, и пересечение любых двух главных идеалов является главным идеалом. Любое кольцо главных идеалов факториально. Дедекиндово кольцо факториально, только если оно-кольцо главных идеалов. Если $S$-мультипликативная система в Ф. к. $A$, то кольцо частных $S^{-1} A$ ф্রакториально. Для кольца Зариского $R$ из факториальности его пополнения $\widehat{R}$ следует факториальность самого $R$.



Подкольцо и факторкольцо Ф. к. не обязаны быть Ф. к. Кольцо многочленов над Ф. к. и кольцо формальных степенных рядов над полем или дискретно нормированным кольцом факториально. Однако кольцо формальных степенных рядов над Ф. к. не обязано быть факториальным.



Область целостности факториальна тогда и только тогда, когда ее мультипликативная полугруппа гауссова (см. Гауссова полугруппа), в связи с әтим Ф.к. наз. также г а у с с о в ы м и кол ь ц а ми.



Лит.: [1] Б у р б а к и Н., Коммутативная алгебра, пер. с франц., М., 1971 .



ФАКТОРИЗАЦИОННАЯ Т ф а к то р и з а ц и и,- теорема теории статистич. оценивания, указывающая необходимое и достаточное условия того, чтобы статистика $T$ была достатотной для семейства вероятностных расп ределений $\left\{P_{\theta}\right\}$.



Пусть $X$-случайный вектор, принимающий значе-



\includegraphics[max width=\textwidth, center]{2023_07_09_b0c21899c65d7aea98ddg-1}

притем семейство вероятностных распределений $\left\{P_{\theta}\right\}$ доминировано нек-рой мерой $\mu$ и пусть



\[

p(x ; \theta)=d P_{\theta}(x) / d \mu, \quad \theta \in \Theta .

\]



Далее, пусть $T=T(X)$-статистика, построенная по вектору наблюдений $X$ и отображающая измеримое пространство $(\mathfrak{X}, \mathscr{B})$ в измеримое пространство $(\mathfrak{A}, \mathcal{A})$. В этих условиях возникает вопрос: когда статистика $T$ является достатотной для семейства $\left\{P_{\theta}\right\}$ ? Отвечая на этот вопрос, Ф. т. утверждает: для того чтобы статистика $T$ была достаточной для семейства $\left\{P_{\theta}\right\}$, допускающего достаточные статистики, необходимо и достаточно, чтобы для любого $\theta \in \Theta$ плотность вероятности $p(x ; \theta)$ можно было факторизовать следующим образом:



\[

p(x ; \theta)=g(x) h(T(x) ; \theta),

\]



где $g(\cdot)$ есть $\mathscr{B}$-измеримая функция на $(\mathfrak{X}, \mathscr{B}), h(\cdot, \theta)$





\end{document}